\section{Rigid-body multirotor dynamics}
\label{sec:rigid-body}

The baseline vehicle model is \texttt{Sim.Vehicles.GenericMultirotor}, a rigid-body multirotor with configurable propulsor geometry. It is intentionally low-fidelity but allocation-free and stable for closed-loop testing.

\subsection{Translational dynamics}
The position derivative is simply
\begin{equation}
  \dot{\vect{p}}_n = \vect{v}_n.
\end{equation}

The total thrust force in body coordinates is formed by summing per-rotor thrust vectors. Each propulsor has a body-frame axis \(\vect{a}_{b,i}\) (\code{QuadrotorParams.rotor\_axis\_body}) with the convention
\begin{equation}
  \vect{F}_{b,i} = -T_i\,\vect{a}_{b,i},
\end{equation}
so for a classic multirotor with thrust along \(-z_b\), \(\vect{a}_{b,i}=(0,0,1)\). The total body force is
\begin{equation}
  \vect{F}_b = \sum_{i=1}^{N} \vect{F}_{b,i}, \qquad \vect{F}_n = \Rbn\,\vect{F}_b.
\end{equation}

A simple linear drag term is applied in NED based on air-relative velocity
\begin{equation}
  \vect{v}_{rel} = \vect{v}_n - \vect{w}_n, \qquad \vect{F}_{drag} = -c_d\,\vect{v}_{rel},
\end{equation}
where \(c_d\) is \code{QuadrotorParams.linear\_drag} and \(\vect{w}_n\) is the held wind sample from the runtime bus.

Finally, gravity is supplied by the environment model as an acceleration vector \(\vect{g}_n\) (positive down in NED). The velocity derivative is
\begin{equation}
  \dot{\vect{v}}_n = \frac{1}{m}\left(\vect{F}_n + \vect{F}_{drag}\right) + \vect{g}_n.
\end{equation}

\subsection{Rotational dynamics with rotor gyroscopics}
Moments from propulsors have two contributions: (i) thrust at a lever arm, and (ii) reaction torque about the propulsor axis. With body-frame rotor position \(\vect{r}_{b,i}\) and thrust \(T_i\),
\begin{equation}
  \vect{\tau} = \sum_{i=1}^{N} \left( \vect{r}_{b,i} \cross \vect{F}_{b,i} + \vect{a}_{b,i}\,\tau_i \right),
\end{equation}
where \(\tau_i\) is the signed shaft torque reported by the propulsion layer. In the current implementation this is the \emph{aerodynamic load torque} \(Q_i\) from the prop model, signed by \texttt{rotor\_dir[i]} (reaction torque on the body). Transient motor spin-up/spin-down effects enter separately through the rotor angular momentum terms \(-\dot{\vect{H}}\) and \(-\vect{\omega}_b \cross \vect{H}\); they are not double-counted inside \(\tau_i\).

The rigid-body angular dynamics include angular damping and the additional angular momentum stored in spinning rotors. Let \(\mat{I}\) be the body inertia, \(\vect{\omega}_b\) the body rates, and define the rotor angular momentum and its derivative
\begin{equation}
  \vect{H} = \sum_{i=1}^{N} J_i\,s_i\,\Omega_i\,\vect{a}_{b,i}, \qquad
  \dot{\vect{H}} = \sum_{i=1}^{N} J_i\,s_i\,\dot{\Omega}_i\,\vect{a}_{b,i},
\end{equation}
where \(J_i\) is the axial inertia of rotor \(i\), \(\Omega_i\) its spin rate (rad/s), and \(s_i=-\texttt{rotor\_dir}[i]\) is the rotor spin sign (opposite the reaction-torque sign). The implemented equation of motion is
\begin{equation}
  \mat{I}\,\dot{\vect{\omega}}_b = \left(\vect{\tau} - \vect{d}\odot\vect{\omega}_b\right)
  - \vect{\omega}_b\cross(\mat{I}\vect{\omega}_b)
  - \vect{\omega}_b\cross\vect{H} - \dot{\vect{H}},
\end{equation}
with elementwise angular damping \(\vect{d}\) from \code{QuadrotorParams.angular\_damping}. Quaternion kinematics use \cref{sec:frames}.

\begin{figure}[ht]
  \centering
  \includegraphics[width=0.75\linewidth]{figs/gyro_coupling_plot.png}
  \caption{Gyro coupling regression for a tilted rotor axis: expected vs simulated \(\dot{\omega}\) components.}
  \label{fig:gyro-coupling}
\end{figure}

\subsection{Defaults and caveats}
\paragraph{Defaults.} \code{QuadrotorParams} provides constructor defaults (used if a spec omits values): mass \SI{1.5}{kg}, diagonal inertia \((0.029125,0.029125,0.055225)\,\si{kg.m^2}\), linear drag \SI{0.05}{N/(m/s)}, and angular damping $(0.02,0.02,0.01)$.

\paragraph{Contact forces.} The rigid-body model itself does not include ground contact. Contact is handled in the coupled plant layer: penalty contact adds an external COM force in the RHS, while constraint contact can use hybrid interval stepping (time-of-impact + impact map + grounded impulses). See \cref{sec:contacts}.

\paragraph{Aerodynamics.} The drag term is deliberately simplistic (linear, isotropic, no lift/induced effects). This is appropriate for controller-in-the-loop regression but not for aerodynamic performance prediction.

\subsection{TOML configuration: airframe geometry and initial state}
\label{sec:airframe-toml}

The rigid-body parameters in \cref{sec:rigid-body} are configured under \code{[airframe]} (\code{AirframeSpec} in \code{src/sim/Aircraft/Spec.jl}) and its nested \code{[airframe.x0]} initial-condition table. A typical multirotor configuration (matching the shipped Iris asset) is:

\begin{lstlisting}[language=toml]
[airframe]
kind = "multirotor"
mass_kg = 1.5
inertia_diag_kgm2 = [0.029125, 0.029125, 0.055225]
# Optional products of inertia Ixy, Ixz, Iyz (defaults to zeros)
inertia_products_kgm2 = [0.0, 0.0, 0.0]

rotor_pos_body_m = [
  [ 0.1515,  0.2450, 0.0],
  [-0.1515, -0.1875, 0.0],
  [ 0.1515, -0.2450, 0.0],
  [-0.1515,  0.1875, 0.0],
]
rotor_axis_body_m = [
  [0.0, 0.0, 1.0],
  [0.0, 0.0, 1.0],
  [0.0, 0.0, 1.0],
  [0.0, 0.0, 1.0],
]

linear_drag = 0.05
angular_damping = [0.02, 0.02, 0.01]

[airframe.x0]
pos_ned    = [0.0, 0.0, 0.0]
vel_ned    = [0.0, 0.0, 0.0]
q_bn       = [1.0, 0.0, 0.0, 0.0]   # or euler_deg / euler_rad
omega_body = [0.0, 0.0, 0.0]
\end{lstlisting}

\paragraph{Parameter mapping to the dynamics.}
The keys above map one-to-one into the symbols used in \cref{sec:rigid-body}:
\begin{itemize}[leftmargin=*,itemsep=2pt]
  \item \code{mass\_kg} \(\rightarrow m\).
  \item \code{inertia\_diag\_kgm2} and \code{inertia\_products\_kgm2} \(\rightarrow\) the symmetric inertia tensor \(\mat{I}\). (Alternatively, \code{inertia\_kgm2} may be supplied as a full symmetric \(3\times 3\) tensor.)
  \item \code{rotor\_pos\_body\_m} \(\rightarrow \vect{r}_{b,i}\) for each rotor.
  \item \code{rotor\_axis\_body\_m} \(\rightarrow \vect{a}_{b,i}\) for each rotor. Axes are normalized during build; if omitted entirely, all axes default to \((0,0,1)\) (thrust along \(-z_b\)).
  \item \code{linear\_drag} \(\rightarrow c_d\) in \(\vect{F}_{drag}=-c_d\,\vect{v}_{rel}\).
  \item \code{angular\_damping} \(\rightarrow \vect{d}\) in \(\vect{d}\odot\vect{\omega}_b\).
\end{itemize}

\paragraph{Initial condition semantics.}
The \code{[airframe.x0]} table sets the initial rigid-body state. Attitude may be specified either as \code{q\_bn} or as Euler angles; in strict mode, only one representation may be supplied.
